\section{Small numerical example}
\label{sec:example}

To instantiate the proposed methodology and demonstrate the origin of the stalling geometry, we define a concrete, two-dimensional scenario heavily motivated by the initial orbit determination problem in aerospace engineering. In such nonlinear dynamical systems, initial Gaussian uncertainties propagating through orbital mechanics rapidly deform into highly non-Gaussian geometries, commonly referred to as ``crescent'' distributions. 

\subsection{Target distribution generation}

We begin by generating a source ensemble of $M$ unperturbed particles from a standard bivariate normal distribution, denoted ${\zeta}^{(i)} \sim \mathcal{N}({0}, {I}_2)$. To simulate the nonlinear orbital shearing effect, we apply a deterministic mapping to generate our target crescent-shaped dataset ${z}^{(i)}$:
\begin{align*}
    z_1^{(i)} &= \zeta_1^{(i)}, \\
    z_2^{(i)} &= \zeta_2^{(i)} + \left(\zeta_1^{(i)}\right)^2.
\end{align*}
The fundamental task of our KR optimal transport map ${S}({z})$ is to invert this nonlinear distortion, thereby returning the highly coupled crescent empirical distribution back to a perfectly independent standard normal form.

\subsection{Formulating the one-dimensional map component}

Let us focus explicitly on computing the first component of the triangular map, $S^1(z_1)$, which depends solely on the first spatial coordinate. We truncate our parameterisation at a maximum Hermite polynomial degree of $D=3$. Consequently, our parameterised map is defined as:
\begin{equation*}
    S^1(z_1) = w_0 \text{He}_0(z_1) + w_1 \text{He}_1(z_1) + w_2 \text{He}_2(z_1) + w_3 \text{He}_3(z_1).
\end{equation*}
Substituting the explicit algebraic forms of the probabilist's Hermite polynomials, we obtain the expression:
\begin{equation*}
    S^1(z_1) = w_0(1) + w_1(z_1) + w_2(z_1^2 - 1) + w_3(z_1^3 - 3z_1).
\end{equation*}
By applying the derivative property from equation \eqref{eq:hermite_derivative}, the derivative of the map with respect to $z_1$ is trivially derived as:
\begin{equation*}
    \partial_1 S^1(z_1) = w_1(1) + 2 w_2 (z_1) + 3 w_3 (z_1^2 - 1).
\end{equation*}

\subsection{Constructing the one-dimensional constraint matrix}

To enforce the KR monotonicity condition across the entire ensemble, we require that $\partial_1 S^1(z_1^{(i)}) \ge \varepsilon$ holds for every particle index $i$. For a single, specific particle located at coordinate $z_1^{(i)}$, this constraint evaluates algebraically to the linear inequality:
\begin{equation*}
    0 \cdot w_0 + 1 \cdot w_1 + 2z_1^{(i)} \cdot w_2 + 3\left((z_1^{(i)})^2 - 1\right) \cdot w_3 \ge \varepsilon.
\end{equation*}
This inequality dictates a single row of our constraint matrix ${A}$. Therefore, the $i$-th row ${A}_{i,:}$ is given by the vector:
\begin{equation*}
    {A}_{i,:} = \left[ 0, \quad 1, \quad 2z_1^{(i)}, \quad 3\left((z_1^{(i)})^2 - 1\right) \right].
\end{equation*}

\subsection{Formulating the two-dimensional map component}

Having Gaussianised the first coordinate, we now extend our formulation to the second component of the triangular map, $S^2(z_1, z_2)$, which depends on both spatial coordinates. We approximate this component using a tensor product of our probabilist's Hermite polynomials up to a maximum degree $D$. The parameterised map is defined as the bivariate expansion:
\begin{equation*}
    S^2(z_1, z_2) = \sum_{j=0}^D \sum_{k=0}^D w_{j,k} \text{He}_j(z_1) \text{He}_k(z_2),
\end{equation*}
where $w_{j,k}$ are the scalar coefficients to be optimised. 

The KR monotonicity condition dictates that this second component must be strictly increasing with respect to its second argument. Consequently, we must ensure that $\partial_2 S^2(z_1, z_2) \ge \varepsilon$. By exploiting the derivative property of the Hermite polynomials specifically along the $z_2$ dimension, the partial derivative simplifies analytically to:
\begin{equation*}
    \partial_2 S^2(z_1, z_2) = \sum_{j=0}^D \sum_{k=1}^D w_{j,k} \text{He}_j(z_1) \left( k \text{He}_{k-1}(z_2) \right).
\end{equation*}

\subsection{Constructing the two-dimensional constraint matrix}

To enforce the monotonicity condition across the entire continuous cloud, we evaluate the partial derivative at every empirical particle $i$. For a specific particle located at $(z_1^{(i)}, z_2^{(i)})$, the constraint reduces to the linear inequality:
\begin{equation*}
    \sum_{j=0}^D \sum_{k=1}^D w_{j,k} \text{He}_j(z_1^{(i)}) \left( k \text{He}_{k-1}(z_2^{(i)}) \right) \ge \varepsilon.
\end{equation*}
By computing the constant terms $\text{He}_j(z_1^{(i)}) (k \text{He}_{k-1}(z_2^{(i)}))$ for all combinations of indices $j$ and $k$, we construct the $i$-th row of our expanded constraint matrix ${A}$ for the two-dimensional problem.